\scanpage{58}
On posera \deleted{enfin}
\[
(5)\qquad V_{ij}=\overline V_{ij}\setminus\dot V_{ij}=\overline V_{ij}-\overline V_{ij}\cap R_{ij}
\]
qui est donc un espace infinitésimal autour de
\[
(6)\qquad V_{ij}\cap X_i=X_i^*.
\]
Je poserai enfin
\[
(7)\qquad\begin{aligned}
V_{ij}^*&=\overline V_{ij}\setminus\bigcup_{\deleted{\unsure{i}}\,\uline{$k$}<j}X_k\cap\overline V_{ij}\\
&=\overset{\deleted{\unsure{-}}}{\overline V}_{ij}\setminus\dot X_j\\
&=V_{ij}\cap X_j^*\subset V_{i,j}.
\end{aligned}
\]
Notons que pour $i=j$, on trouve
\[
(8)\qquad\overline V_{ii}=X_i,\quad\dot V_{i,i}=\dot X_i,\quad\uline{$V_{i,i}=V_{i,i}^*=X_i^*$}
\]
Nos ``pierres'' de construction, dans une première étape, vont être les $V_{ij}^*$, $V_{ij}$ avec les inclusions évidentes
\diagram{(9)\qquad\begin{tikzcd}[column sep=large,row sep=small]
& X_i^*\arrow[d,"\text{``âme''}"]\\
\unsure{V_{ij}^*|X_{ij}^*}\arrow[r]\arrow[d] & V_{i,j}\\
X_j^* &
\end{tikzcd}}
,
ces pierres \unsure{comprenant|comprennent} \inserted{les} $X_i^*$ comme
\mspage{57}
